\documentclass[11pt,a4paper]{article}

\usepackage[T1]{fontenc}
\usepackage[utf8]{inputenc}
\usepackage{lmodern}
\usepackage{microtype}
\usepackage{geometry}
\geometry{margin=2.5cm}

\usepackage[hidelinks,
  colorlinks=true,
  linkcolor=blue!70!black,
  urlcolor=blue!70!black,
  citecolor=blue!70!black,
  pdftitle={ChemEditor Engineering Guide},
  pdfauthor={ChemEditor},
]{hyperref}

\usepackage{amsmath}
\usepackage{booktabs}
\usepackage{tabularx}
\usepackage{listings}
\usepackage{xcolor}
\usepackage{parskip}
\usepackage{titlesec}
\usepackage{enumitem}
\usepackage{fancyhdr}
\usepackage{graphicx}

% ── Code listing style ──────────────────────────────────────────────────────
\definecolor{codebg}{HTML}{F5F5F5}
\definecolor{codeframe}{HTML}{CCCCCC}
\definecolor{kwcolor}{HTML}{0000CC}
\definecolor{cmtcolor}{HTML}{007700}
\definecolor{strcolor}{HTML}{CC0000}

\lstset{
  backgroundcolor=\color{codebg},
  frame=single,
  framerule=0.4pt,
  rulecolor=\color{codeframe},
  basicstyle=\ttfamily\small,
  keywordstyle=\color{kwcolor}\bfseries,
  commentstyle=\color{cmtcolor}\itshape,
  stringstyle=\color{strcolor},
  breaklines=true,
  breakatwhitespace=true,
  tabsize=2,
  showstringspaces=false,
  captionpos=b,
  upquote=true,
}

% ── Custom language definitions ───────────────────────────────────────────
\lstdefinelanguage{TypeScript}{
  keywords={const,let,var,function,return,if,else,for,while,class,interface,
            type,import,export,from,default,async,await,new,this,null,
            undefined,true,false,extends,implements,typeof,instanceof,in,
            of,try,catch,throw,void,any,string,number,boolean,never,
            unknown,readonly,as,keyof,Record,Set,Map,Promise,Array,
            useCallback,useEffect,useMemo,useRef,useState,forwardRef,
            useImperativeHandle,lazy,Suspense},
  sensitive=true,
  comment=[l]{//},
  morecomment=[s]{/*}{*/},
  morestring=[b]',
  morestring=[b]",
  morestring=[b]`,
}

\lstdefinelanguage{Rust}{
  keywords={fn,let,mut,const,use,pub,mod,impl,struct,enum,trait,for,in,
            if,else,while,loop,match,return,self,Self,super,crate,async,
            await,move,ref,where,type,dyn,unsafe,extern,static,Some,None,
            Ok,Err,true,false,String,Vec,Option,Result},
  sensitive=true,
  comment=[l]{//},
  morecomment=[s]{/*}{*/},
  morestring=[b]",
}

% ── Section formatting ───────────────────────────────────────────────────
\titleformat{\section}{\large\bfseries}{}{0em}{}[\titlerule]
\titleformat{\subsection}{\normalsize\bfseries}{}{0em}{}
\titleformat{\subsubsection}{\normalsize\itshape\bfseries}{}{0em}{}

% ── Header / footer ──────────────────────────────────────────────────────
\pagestyle{fancy}
\fancyhf{}
\fancyhead[L]{\small\textit{ChemEditor --- Engineering Guide}}
\fancyhead[R]{\small\textit{\today}}
\fancyfoot[C]{\thepage}
\renewcommand{\headrulewidth}{0.4pt}

% ── Convenience macros ───────────────────────────────────────────────────
% Inline code with underscores and special chars
\newcommand{\ci}[1]{\lstinline|#1|}

% ─────────────────────────────────────────────────────────────────────────────
\begin{document}

\begin{center}
  {\LARGE\bfseries ChemEditor --- Engineering Guide}\\[0.6em]
  {\large A pedagogical deep-dive for people fluent in Python/Julia/C++ who\\
  have never used Tauri, React, Rust, or TypeScript}\\[1em]
  \textit{\today}
\end{center}

\noindent\textbf{How to use this document.}
This guide is the long-form conceptual architecture reference.
For current setup instructions, package scripts, ChemDraw fidelity workflows, and NMR data
maintenance, start with the Markdown docs under \texttt{development/}, especially
\texttt{development/README.md} and \texttt{development/getting-started.md}.

\tableofcontents
\newpage

% ─────────────────────────────────────────────────────────────────────────────
\section{What Is This Project?}
\label{sec:what}

ChemEditor is a \textbf{native desktop molecule drawing editor} --- think a lightweight
ChemDraw --- built with a web UI inside a Tauri shell. You draw chemical structures
with your mouse; the app converts them in real time to molecular representations,
supports ChemDraw-compatible CDXML import and export, and can generate and inspect 3D
conformers. The current operational details of the build, sidecar, and fidelity harness
live in the Markdown docs under \texttt{development/}; this document focuses on the
architectural ideas behind the implementation.

\subsection{Why These Technology Choices?}

\begin{table}[h!]
\centering
\begin{tabularx}{\textwidth}{@{} l l X @{}}
\toprule
\textbf{Need} & \textbf{Choice} & \textbf{Why} \\
\midrule
Native desktop app & Tauri & Small binary ($\sim$10\,MB); uses the OS's built-in web
  renderer --- no bundled Chromium like Electron \\
Interactive drawing canvas & React + react-konva & Konva is an HTML5 Canvas library
  wrapped in React; lets you describe canvas geometry declaratively \\
Real-time SMILES / 2D structure & RDKit WASM & RDKit compiled to WebAssembly runs
  in-process --- no network round-trips, no subprocess latency \\
3D conformer generation & Python sidecar (RDKit) & Heavy ETKDGv3 computation; Python's
  RDKit is the gold standard; \texttt{uv} manages the environment cleanly \\
3D viewer & Custom Three.js viewer & Direct control over camera, lighting,
  measurements, conformer alignment, and future volumetric rendering \\
NMR shift prediction & HOSE lookup DB + rule corrections & Offline HOSE database
  generated from NMRShiftDB2 (ODbL) and queried in the Python sidecar \\
ChemDraw compatibility & Custom CDXML writer & CDXML is XML; we write it ourselves
  to control atom/bond format and coordinate system precisely \\
\bottomrule
\end{tabularx}
\end{table}

% ─────────────────────────────────────────────────────────────────────────────
\section{The Technology Stack At a Glance}
\label{sec:stack}

If you know Python, Julia, or C++, each piece of this stack has a near-perfect analogy:

\begin{table}[h!]
\centering
\begin{tabularx}{\textwidth}{@{} l l X @{}}
\toprule
\textbf{Technology} & \textbf{What it is} & \textbf{Analogy} \\
\midrule
\textbf{Tauri} & Framework for native desktop apps with a web UI
  & Python + PyQt, but the UI is an HTML page and the backend is Rust \\
\textbf{React} & UI library that re-renders when data changes
  & Jinja2 templates that re-run themselves whenever a variable changes \\
\textbf{TypeScript} & JavaScript with type annotations
  & Python with \texttt{typing} --- but types are enforced at compile time \\
\textbf{Rust} & Systems language for the Tauri backend
  & C++ without undefined behaviour, with ownership enforced by the compiler \\
\textbf{Vite} & JavaScript build tool and dev server
  & CMake for JavaScript; bundles source files into optimised output \\
\textbf{pnpm} & JavaScript package manager
  & \texttt{pip} --- installs packages from a registry into \texttt{node\_modules/} \\
\textbf{Cargo} & Rust package manager and build system
  & \texttt{pip} + CMake combined; reads \texttt{Cargo.toml}, compiles everything \\
\textbf{Zustand} & Global state library for React
  & A Python class where every attribute setter notifies all subscribers \\
\textbf{react-konva} & React bindings for the Konva Canvas 2D library
  & \texttt{matplotlib} but interactive --- describes shapes on an HTML5 canvas \\
\textbf{RDKit WASM} & RDKit compiled to WebAssembly
  & RDKit, but running inside the browser sandbox without a server \\
\textbf{Three.js viewer} & Custom WebGL small-molecule viewer
  & A focused in-app molecular renderer built directly on top of \texttt{three.js} \\
\textbf{uv} & Fast Python package and project manager
  & \texttt{pip} + \texttt{venv} + \texttt{pyproject.toml} all in one tool \\
\textbf{PyInstaller} & Packages Python scripts into standalone executables
  & Like \texttt{cx\_Freeze} --- produces a single binary with Python bundled inside \\
\textbf{HOSE lookup DB} & Precomputed NMR environment database
  & Canonical local-environment lookup table for fast laptop-scale prediction without a
    neural inference runtime \\
\bottomrule
\end{tabularx}
\end{table}

% ─────────────────────────────────────────────────────────────────────────────
\section{Tauri In Depth}
\label{sec:tauri}

\subsection{What Tauri Is}

Tauri lets you build a native desktop application where the UI is rendered by the
operating system's built-in web engine (WebKit on Linux/macOS, WebView2 on Windows)
and the backend logic runs as a compiled Rust binary.

The result is a \texttt{.deb}/\texttt{.rpm} (on Linux) that weighs $\sim$10\,MB instead
of Electron's $\sim$200\,MB, because you are not bundling a copy of Chrome. The web
engine is already on every modern OS; Tauri just uses it.

\subsection{The Two-Process Model}

\begin{lstlisting}[language={}]
+------------------------------------------+
|  WebView (your React app)                |
|  Runs JavaScript, renders HTML/CSS       |
|  Has NO file system access by default    |
|  Lives in a security sandbox             |
+---------------------+--------------------+
                      |  IPC (invoke / emit)
                      |  serialised JSON
+---------------------v--------------------+
|  Rust backend (src-tauri/src/)           |
|  Has full OS access                      |
|  Manages plugins: fs, dialog, shell      |
|  Runs the Python sidecar as a subprocess |
+------------------------------------------+
\end{lstlisting}

Think of it like a Python web app where the browser is the frontend and a
Flask/FastAPI server is the backend --- except everything runs in the same desktop
process with no HTTP socket. Communication is via IPC, and the ``server'' is a
compiled Rust binary rather than a Python process.

\subsection{IPC: Calling Rust from JavaScript}

The core pattern is: annotate a Rust function with \lstinline|#[tauri::command]|,
then call it from TypeScript with \lstinline|invoke()|.

\textbf{Rust side}:
\begin{lstlisting}[language=Rust]
#[tauri::command]
async fn generate_conformers(
    app: tauri::AppHandle,
    smiles: String,
    force_field: Option<String>,
    max_conformers: Option<i32>,
) -> Result<serde_json::Value, String> {
    // ... do work ...
}
\end{lstlisting}

\textbf{TypeScript side}:
\begin{lstlisting}[language=TypeScript]
import { invoke } from '@tauri-apps/api/core';

const result = await invoke('generate_conformers', {
  smiles: 'CCO',
  forceField: 'MMFF94s',
  maxConformers: 5,
});
\end{lstlisting}

Tauri automatically converts the function name from \texttt{snake\_case} to
\texttt{camelCase}. Arguments are serialised to JSON at the TypeScript boundary and
deserialised back to Rust types automatically via the \texttt{serde} crate. You never
write serialisation code.

The command must be registered at startup; \lstinline|generate_handler!| is a Rust
macro that builds the dispatch table at compile time:

\begin{lstlisting}[language=Rust]
// src-tauri/src/lib.rs
tauri::Builder::default()
    .invoke_handler(tauri::generate_handler![
        generate_conformers,
        parse_chemical_file,
        export_chemical_file
    ])
\end{lstlisting}

\subsection{Tauri Plugins}

Tauri itself only handles IPC. Additional OS capabilities come from plugins
initialised in the same \texttt{Builder} chain:

\begin{table}[h!]
\centering
\begin{tabularx}{\textwidth}{@{} l l X @{}}
\toprule
\textbf{Plugin} & \textbf{What it gives us} & \textbf{Analogous Python call} \\
\midrule
\texttt{tauri-plugin-fs}     & Read/write files on disk & \texttt{open()} / \texttt{pathlib} \\
\texttt{tauri-plugin-dialog} & Native OS open/save file dialogs & \texttt{tkinter.filedialog} \\
\texttt{tauri-plugin-shell}  & Spawn subprocesses; manage sidecars & \texttt{subprocess.Popen} \\
\bottomrule
\end{tabularx}
\end{table}

\subsection{Configuration: tauri.conf.json}

\begin{lstlisting}[language={}]
{
  "app": {
    "windows": [{ "title": "ChemEditor", "width": 800, "height": 600,
                  "decorations": false }],
    "security": { "csp": null }
  },
  "bundle": {
    "targets": ["deb", "rpm"],
    "externalBin": ["bin/chem-engine"]
  }
}
\end{lstlisting}

Key fields:
\begin{itemize}[leftmargin=1.5em]
  \item \lstinline|"decorations": false| --- removes the OS title bar. We draw our
        own minimize/maximize/close buttons in \texttt{App.tsx}.
  \item \lstinline|"csp": null| --- disables the Content Security Policy, letting
        RDKit WASM and the viewer's dynamically loaded assets run without additional
        CSP plumbing.
  \item \lstinline|"externalBin"| --- tells Tauri to bundle the named binary
        alongside the app. The actual file on disk must be named with the full Rust
        target triple (explained in the Build System section).
\end{itemize}

% ─────────────────────────────────────────────────────────────────────────────
\section{Rust Crash Course}
\label{sec:rust}

You don't need to write Rust to work on this project, but you need to read
\texttt{src-tauri/src/lib.rs}. Here is the minimum you need.

\subsection{Syntax Basics}

\begin{lstlisting}[language=Rust]
// fn = function (like def in Python, or just a function in Julia)
fn add(x: i32, y: i32) -> i32 {
    x + y   // No semicolon = implicit return. Identical to Julia's convention.
}

// async fn = coroutine (like Python's async def)
async fn fetch_data() -> String {
    some_async_call().await  // .await comes after the expression, not before
}

// let = variable binding (like Python's = or Julia's assignment)
let x: i32 = 42;   // type annotation after colon, like Python's x: int = 42
let y = 42;         // type is inferred -- Rust has full type inference like Julia
\end{lstlisting}

\subsection{Attributes (Decorators) and Macros}

Rust has no decorators, but it has \textbf{attributes} --- compile-time metadata
that tools like Tauri read to generate boilerplate. They look like Python decorators:

\begin{lstlisting}[language=Rust]
#[tauri::command]   // like @tauri.command -- marks this fn as callable-from-JS
async fn my_command(input: String) -> Result<String, String> { ... }

#[cfg_attr(mobile, tauri::mobile_entry_point)]  // conditional attribute
pub fn run() { ... }
\end{lstlisting}

The \lstinline|#[...]| syntax is an \textbf{attribute}. Macros invoked with \lstinline|!|
are a separate thing:
\begin{lstlisting}[language=Rust]
tauri::generate_handler![cmd_a, cmd_b]  // the ! makes this a macro call
println!("hello {}", name)              // also a macro -- Rust's print
\end{lstlisting}

Think of macros as code-generation functions that run at compile time, before the
compiler processes your code.

\subsection{Result: Error Handling Without Exceptions}

Python uses exceptions. Rust uses \lstinline|Result<T, E>|, which forces you to handle
errors explicitly as a value:

\begin{lstlisting}[language=Rust]
// Result<Value, String> is like Union[dict, str] in Python:
// either a successful value (Ok) or an error message (Err)

Ok(my_value)             // equivalent to: return my_value
Err("something failed")  // equivalent to: raise Exception("something failed")
\end{lstlisting}

Tauri serialises \lstinline|Ok(v)| as the JSON response, and \lstinline|Err(s)| as an
error that becomes a JavaScript exception on the calling side.

The \lstinline|?| operator is syntactic sugar for ``return Err if this failed,
otherwise unwrap the Ok'':
\begin{lstlisting}[language=Rust]
let output = some_operation()?;
// Equivalent to:
// let output = match some_operation() {
//     Ok(v)  => v,
//     Err(e) => return Err(e.into()),
// };
\end{lstlisting}

\subsection{Option}

\lstinline|Option<T>| means ``might not be present'':
\begin{lstlisting}[language=Rust]
Option<String>   // Python: Optional[str]
Some("hello")    // Python: "hello"
None             // Python: None
value.unwrap_or_else(|| "default".to_string())  // Python: value or "default"
\end{lstlisting}

\subsection{serde\_json::Value}

This is Rust's equivalent of Python's \texttt{Any}-typed JSON:
\begin{lstlisting}[language=Rust]
serde_json::Value  // Python: dict | list | str | int | float | bool | None
\end{lstlisting}
You can deserialise arbitrary JSON into it and the type system is satisfied.

\subsection{String vs \&str}

\begin{itemize}[leftmargin=1.5em]
  \item \lstinline|String| --- owned, heap-allocated string (like Python's \texttt{str})
  \item \lstinline|&str| --- borrowed reference to a string slice (like C++'s
        \texttt{std::string\_view})
\end{itemize}

For Tauri command parameters, always use \lstinline|String| --- Tauri's deserialiser
needs owned values. String literals like \lstinline|"hello"| in code are \lstinline|&str|.

\subsection{Closures}

\begin{lstlisting}[language=Rust]
// Rust closure                   // Python lambda
let sq = |x: i32| x * x;        // sq = lambda x: x * x
let add = |a, b| a + b;          // add = lambda a, b: a + b
\end{lstlisting}

The \lstinline||| delimiters hold the parameter list; the body follows. Multi-line
closures use \lstinline|{}|:
\begin{lstlisting}[language=Rust]
let result = vec![1, 2, 3].iter().map(|x| x * 2).collect::<Vec<_>>();
// Python: result = [x * 2 for x in [1, 2, 3]]
\end{lstlisting}

\subsection{lib.rs Walkthrough}

\texttt{src-tauri/src/lib.rs} defines three public commands and one private helper.

\textbf{\lstinline|generate_conformers|} --- Takes a SMILES string, a force field
name, and a max-conformers count. Calls \lstinline|run_engine| with the
\lstinline|conformer| subcommand. Returns JSON
\lstinline|{conformers: [...], bonds: [...]}|.

\textbf{\lstinline|parse_chemical_file|} --- Takes file content as a string and a
format specifier (\lstinline|"mol"|, \lstinline|"pdb"|, \lstinline|"xyz"|, etc.).
Passes the content via stdin to the sidecar's \lstinline|parse| subcommand.
Returns \lstinline|{molblock, smiles}|.

\textbf{\lstinline|export_chemical_file|} --- Takes a V2000 Molblock string and a
format. For RXN export also takes an optional JSON string of arrow data.
Returns \lstinline|{content: "..."}| with the formatted file text.

\textbf{\lstinline|run_engine|} --- The shared helper. It tries two paths:

\textit{Production path:} \lstinline|app.shell().sidecar("chem-engine")| asks Tauri to
locate the bundled binary named \lstinline|chem-engine-<target-triple>|. This binary
is a self-contained PyInstaller executable.

\textit{Development path:} If \lstinline|sidecar()| fails (because the binary hasn't
been built yet), the code falls back to:
\begin{lstlisting}[language=Rust]
std::process::Command::new("uv")
    .arg("run").arg("--project").arg(&project_path)
    .arg("python").arg(&script_path)
    .args(args)
\end{lstlisting}
This calls \lstinline|uv run python chem-engine.py| directly --- no binary required
during development.

After spawning, the helper optionally writes \lstinline|stdin_content| to the
process's stdin, then collects all stdout bytes. On exit it parses stdout as JSON
using \lstinline|serde_json::from_slice|. If parsing fails or the process exited
non-zero, it returns an \lstinline|Err|.

\textbf{\texttt{src-tauri/src/main.rs}} contains:
\begin{lstlisting}[language=Rust]
#![cfg_attr(not(debug_assertions), windows_subsystem = "windows")]

fn main() {
    chem_editor_lib::run()
}
\end{lstlisting}

The \lstinline|windows_subsystem = "windows"| attribute suppresses the console
window that would otherwise appear behind the main window in production builds on
Windows. \lstinline|cfg_attr| means ``only apply this attribute when
\lstinline|debug_assertions| are off'' --- i.e., release builds only.

% ─────────────────────────────────────────────────────────────────────────────
\section{TypeScript Crash Course}
\label{sec:ts}

\subsection{TypeScript = JavaScript + Type Annotations}

If you know Python's \texttt{typing} module, TypeScript's type system will feel
familiar --- except the types are enforced at compile time (like C++), not stripped
at runtime (like Python's type hints are by default).

\begin{lstlisting}[language=TypeScript]
// TypeScript                    // Python equivalent
const x: number = 42;           // x: int = 42
let name: string = "caffeine";  // name: str = "caffeine"

function add(a: number, b: number): number {
  return a + b;
}
// Python: def add(a: int, b: int) -> int: return a + b
\end{lstlisting}

\subsection{interface $\approx$ TypedDict}

\begin{lstlisting}[language=TypeScript]
// TypeScript
interface Atom {
  id: string;
  x: number;
  y: number;
  element: string;
  charge?: number;   // ? means optional -- like Optional[int] in Python
}

// Python equivalent
from typing import TypedDict, Optional
class Atom(TypedDict):
    id: str
    x: float
    y: float
    element: str
    charge: Optional[int]
\end{lstlisting}

\subsection{Union Types and Literal Types}

\begin{lstlisting}[language=TypeScript]
type ArrowType = 'reaction' | 'equilibrium' | 'retrosynthetic' | 'curved';
// Python: ArrowType = Literal['reaction', 'equilibrium', 'retrosynthetic', 'curved']

type Tool = 'select' | 'atom' | 'bond' | 'eraser' | 'fragment' | 'arrow';
\end{lstlisting}

\subsection{const and let}

\begin{lstlisting}[language=TypeScript]
const PI = 3.14159;  // Immutable binding -- like Julia's const.
                     // Note: const objects can have their fields mutated;
                     // only the binding itself is immutable.
let x = 5;           // Mutable -- like a normal variable.
// var also exists but is deprecated; never use it in modern TypeScript.
\end{lstlisting}

\subsection{Arrow Functions}

\begin{lstlisting}[language=TypeScript]
const square = (x: number): number => x * x;
// Python: square = lambda x: x * x

const add = (a: number, b: number): number => {
  return a + b;
};

// Without return type annotation (inferred):
const greet = (name: string) => `Hello, ${name}`;
// Python: f"Hello, {name}"
\end{lstlisting}

\subsection{Modules}

\begin{lstlisting}[language=TypeScript]
// Named imports (most common)
import { Atom, Bond } from '../types/chemistry';
// Python: from chemistry import Atom, Bond

// Default import
import React from 'react';

// Named export
export function graphToMolblock(...) { ... }
export const MOLBLOCK_SCALE = 30;

// Type-only import (erased at runtime, used only for type checking)
import type { Atom } from '../types/chemistry';
\end{lstlisting}

\subsection{async/await}

Identical concept to Python's asyncio:
\begin{lstlisting}[language=TypeScript]
// TypeScript
const result = await invoke('generate_conformers', { smiles });

// Python
// result = await generate_conformers(smiles)
\end{lstlisting}

\subsection{Generics}

\begin{lstlisting}[language=TypeScript]
Set<string>              // Python: set[str]
Map<string, number>      // Python: dict[str, int]
Promise<Value>           // Python: Coroutine[..., Value]
Array<Atom>              // Python: list[Atom]
Record<string, number>   // Python: dict[str, int]
\end{lstlisting}

\subsection{Nullish Coalescing and Optional Chaining}

These appear frequently in the codebase:
\begin{lstlisting}[language=TypeScript]
const smi = mol?.get_smiles() ?? '';
// Equivalent to:
// const smi = (mol !== null && mol !== undefined) ? mol.get_smiles() : '';
// Python: smi = mol.get_smiles() if mol is not None else ''
\end{lstlisting}

\lstinline|?.| short-circuits and returns \lstinline|undefined| if the left side is
\lstinline|null| or \lstinline|undefined|, rather than throwing. \lstinline|??| returns
the right side only if the left is \lstinline|null| or \lstinline|undefined| (unlike
\lstinline||||, which also fires for \lstinline|0|, \lstinline|""|, \lstinline|false|).

\subsection{Destructuring}

\begin{lstlisting}[language=TypeScript]
// Object destructuring -- pull named fields out
const { atoms, bonds } = useStore.getState();
// Python: state = useStore.getState(); atoms = state.atoms; bonds = state.bonds

// Array destructuring -- pull by position
const [count, setCount] = useState(0);

// With renaming
const { x: atomX, y: atomY } = atom;
// Python: atomX = atom.x; atomY = atom.y
\end{lstlisting}

\subsection{Template Literals}

\begin{lstlisting}[language=TypeScript]
const msg = `Target: ${triple}`;
// Python: f"Target: {triple}"
\end{lstlisting}

% ─────────────────────────────────────────────────────────────────────────────
\section{React Crash Course}
\label{sec:react}

\subsection{What React Does}

React keeps a UI in sync with data. You describe \textit{what the UI should look like
given the current data}, and React figures out the minimum DOM changes needed to make
it so. You never manually update the UI --- you update the data, and the UI follows.

This is like writing a Julia Makie \lstinline|@lift| expression: you say ``this plot
depends on these observables'', and Makie re-renders automatically when they change.
React is the same idea, but for HTML.

\subsection{Components}

A component is just a function that returns UI:

\begin{lstlisting}[language=TypeScript]
function MoleculePreview({ smiles }: { smiles: string }) {
  return <div>{smiles}</div>;
}
// The argument is called "props" and is a plain object.
// { smiles } is destructuring -- pulling the smiles field out.
\end{lstlisting}

\subsection{JSX}

JSX looks like HTML but it is TypeScript:
\begin{lstlisting}[language=TypeScript]
<Circle x={10} y={20} radius={5} fill="white" />
// Compiles to: React.createElement(Circle, { x: 10, y: 20, radius: 5, fill: "white" })
\end{lstlisting}

Curly braces \lstinline|{}| inside JSX mean ``evaluate this TypeScript expression here'':
\begin{lstlisting}[language=TypeScript]
<div style={{ color: isDarkMode ? '#fff' : '#000' }}>
  {atoms.length} atoms
</div>
\end{lstlisting}

Note the double braces: the outer \lstinline|{}| means ``JS expression'' and the inner
\lstinline|{}| is an object literal.

\subsection{useState: Reactive Variables}

\begin{lstlisting}[language=TypeScript]
const [count, setCount] = useState(0);
// count    -- the current value (read-only inside the component)
// setCount -- function to update it (causes React to re-render)
setCount(count + 1);
setCount(prev => prev + 1);  // functional update -- preferred when new value
                              // depends on old
\end{lstlisting}

This is like a Julia \texttt{Observable}: changing it triggers any code that depends
on it. The key rule is: \textbf{you must call the setter function to change state ---
you cannot do \lstinline|count++|}. React won't know the variable changed unless you
tell it via the setter.

\subsection{useEffect: Side Effects on State Changes}

\begin{lstlisting}[language=TypeScript]
useEffect(() => {
  // Runs whenever `smiles` changes
  fetchConformers(smiles).then(setData);

  // Optional cleanup: runs before the next effect AND on unmount
  return () => abortController.abort();
}, [smiles]);  // dependency array
\end{lstlisting}

Think of it as a slot/signal connection in Qt, or a \lstinline|Makie.on(observable)|
callback. The dependency array controls when it fires:
\begin{itemize}[leftmargin=1.5em]
  \item \lstinline|[smiles]| --- re-runs when \lstinline|smiles| changes
  \item \lstinline|[]| --- runs once on mount (component creation)
  \item omitted --- runs after every render (almost always wrong)
\end{itemize}

\subsection{useMemo: Memoized Computation}

\begin{lstlisting}[language=TypeScript]
const molblock = useMemo(
  () => graphToMolblock(atoms, bonds, width, height),
  [atoms, bonds, width, height]
);
\end{lstlisting}

This is \texttt{functools.lru\_cache} for React: only recomputes when the listed
dependencies change. Without \lstinline|useMemo|, the function would re-run on every
render even if the inputs hadn't changed.

\subsection{Custom Hooks}

Any function starting with \lstinline|use| that calls other hooks is a ``custom hook''
--- reusable stateful logic:

\begin{lstlisting}[language=TypeScript]
function useBondTool() {
  const [isDragging, setIsDragging] = useState(false);
  const onMouseDown = useCallback((pos) => { ... }, []);
  return { isDragging, onMouseDown };
}

// Usage in ChemCanvas:
const { isDragging, onMouseDown } = useBondTool();
\end{lstlisting}

This is like extracting a Python mixin or a Julia module --- packaging stateful logic
so it can be reused or tested independently.

\subsection{useRef: Mutable Values That Don't Trigger Re-render}

\begin{lstlisting}[language=TypeScript]
const stageRef = useRef<Stage | null>(null);
stageRef.current = new Stage(...);   // mutate freely
\end{lstlisting}

\lstinline|useRef| holds a mutable value in \lstinline|.current| that persists across
renders but does \textbf{not} cause a re-render when changed. Use it for DOM references
and for values that callbacks need to read without stale closure issues (covered next).

\subsection{The Stale Closure Problem --- and Why Every Tool Hook Uses Dual State + Ref}
\label{sec:staleclosure}

This is one of the most important React concepts for understanding the tool hooks.

In JavaScript, closures capture variables \textbf{by value at the time the function is
created}. When you create a \lstinline|useCallback| with \lstinline|[]|
(no dependencies), that callback is created once and never re-created. This means it
captures the initial values of any variables it references --- and if those variables
change later, the callback still sees the old values.

Example of the problem:
\begin{lstlisting}[language=TypeScript]
const [count, setCount] = useState(0);

// This callback is created once because deps = []
const onMouseDown = useCallback(() => {
  console.log(count);  // ALWAYS prints 0, even after setCount(5)!
}, []);               // count is "stale" inside here
\end{lstlisting}

The standard fix is to keep a \lstinline|useRef| in sync with the state, and read
from the ref inside callbacks:
\begin{lstlisting}[language=TypeScript]
const [count, setCount] = useState(0);
const countRef = useRef(0);

// Setter that keeps both in sync:
const setCountBoth = (v: number) => { countRef.current = v; setCount(v); };

// Now callbacks can read the live value via ref:
const onMouseDown = useCallback(() => {
  console.log(countRef.current);  // Always sees the current value
}, []);                            // No stale closure
\end{lstlisting}

\textbf{This pattern appears in every tool hook.} In \lstinline|useBondTool|:
\begin{lstlisting}[language=TypeScript]
const [activeAtomId, setActiveAtomId] = useState<string | null>(null);
const activeAtomIdRef = useRef<string | null>(null);

// Single setter that syncs both:
const setActive = (id: string | null) => {
  activeAtomIdRef.current = id;   // ref: for callbacks to read
  setActiveAtomId(id);            // state: to trigger re-render
};
\end{lstlisting}

The \lstinline|useState| variable triggers re-renders (so the canvas updates). The
\lstinline|useRef| variable is read inside \lstinline|onMouseMove|/\lstinline|onMouseUp|
callbacks without stale closures. Both must be updated together whenever the value
changes.

\subsection{forwardRef and useImperativeHandle}

React components normally communicate via props (parent passes data down) and callbacks
(child calls parent's functions). But sometimes a parent needs to call a method on a
child imperatively --- for example, ``save the canvas to disk right now.''

React's answer is \lstinline|forwardRef| + \lstinline|useImperativeHandle|.

\textbf{Step 1} --- Define what methods the parent can call:
\begin{lstlisting}[language=TypeScript]
export interface ChemCanvasRef {
  saveNative:    () => Promise<void>;
  exportPNG:     () => Promise<void>;
  getMolblock:   () => string;
  loadFromSmiles:(smiles: string) => void;
  cleanUp:       () => void;
}
\end{lstlisting}

\textbf{Step 2} --- Wrap the component in \lstinline|forwardRef|, which lets it
receive a \lstinline|ref| from its parent:
\begin{lstlisting}[language=TypeScript]
export const ChemCanvas = forwardRef<ChemCanvasRef, Props>(({ width, height }, ref) => {
  //                                                                         ^^^
  //                         The ref comes in as the second argument

  useImperativeHandle(ref, () => ({
    saveNative:  async () => { /* ... */ },
    getMolblock: () => graphToMolblock(atoms, bonds, width, height),
    // ...
  }));

  return <Stage>...</Stage>;
});
\end{lstlisting}

\textbf{Step 3} --- The parent creates a ref and passes it to the component:
\begin{lstlisting}[language=TypeScript]
// In App.tsx:
const canvasRef = useRef<ChemCanvasRef>(null);

<ChemCanvas ref={canvasRef} width={w} height={h} />

// Now the parent can call methods imperatively:
await canvasRef.current?.saveNative();
const molblock = canvasRef.current?.getMolblock();
\end{lstlisting}

Both \lstinline|ChemCanvas| and \lstinline|Molecule3DThreePanel| use this pattern.
\lstinline|Molecule3DThreePanel| exposes \lstinline|captureViewer()| so the parent can
export a high-resolution PNG from the active Three.js renderer.

\subsection{lazy() and Suspense --- Deferred Loading}

The 3D viewer bundle is sizeable. We don't want it loaded until the user actually
opens the viewer, so React's \lstinline|lazy()| and \lstinline|Suspense| defer it:

\begin{lstlisting}[language=TypeScript]
// In App.tsx -- lazy import:
const Molecule3DThreePanel = lazy(() => import("./components/Molecule3DThreePanel"));
//                              ^^^^^^^^^^^^^^^^^^^^^^^^^^^^^^^^^^^^^^^^^^^^^^^^^^^^^
//                              This file is not loaded until the viewer is first rendered.

// In JSX -- Suspense wraps the lazy component:
<Suspense fallback={null}>
  <Molecule3DThreePanel ref={molecule3DRef} smiles={viewerSmiles} ... />
</Suspense>
\end{lstlisting}

When \lstinline|Molecule3DThreePanel| first needs to render, React dynamically downloads its
JavaScript bundle, executes it, then renders the component. While loading, it shows
the \lstinline|fallback| (here \lstinline|null| --- nothing).

\subsection{Zustand: Global State}

React's \lstinline|useState| is local to one component. For state that multiple
components need, we use Zustand. It behaves like a global mutable object with
change notifications:

\begin{lstlisting}[language=TypeScript]
// Read individual fields (component re-renders only when this field changes):
const tool  = useStore(s => s.tool);
const atoms = useStore(s => s.atoms);

// Read multiple fields:
const { atoms, bonds } = useStore(s => ({ atoms: s.atoms, bonds: s.bonds }));

// Call actions from anywhere, even outside React components:
useStore.getState().setTool('bond');
useStore.getState().pushToHistory({ atoms, bonds, arrows });
\end{lstlisting}

The selector function \lstinline|s => s.tool| lets Zustand only notify this component
when \lstinline|tool| changes, not when any other field changes. Without selectors,
every store update would re-render every subscriber.

% ─────────────────────────────────────────────────────────────────────────────
\section{Project Architecture Deep Dive}
\label{sec:arch}

\subsection{7a. The Data Model (src/types/chemistry.ts)}

The entire canvas is represented as a chemical graph. This file is tiny but
foundational --- every other file depends on these types:

\begin{lstlisting}[language=TypeScript]
interface Atom {
  id: string;        // UUID -- stable identity across renders, undo/redo, re-parses
  x: number;         // canvas pixel position; origin = top-left corner
  y: number;
  element: string;   // "C", "N", "O", "Ph", "OMe", "tBu", etc.
                     // Can be a shorthand label, not just an element symbol
  label?: string;    // display override (kept for CDXML compatibility)
  charge?: number;   // formal charge (+1, -1, +2, etc.)
  electrons?: number; // number of lone pair electrons for display (0, 2, 4)
}

interface Bond {
  id: string;
  from: string;    // atom ID of first endpoint
  to: string;      // atom ID of second endpoint
  order: number;   // 1=single, 2=double, 3=triple
  stereo?: number; // 0=none, 1=wedge (bold), 6=hash (dashed)
                   // These are V2000 MDL format codes -- see section 7e
}

interface Arrow {
  id: string;
  type: ArrowType;  // 'reaction' | 'equilibrium' | 'retrosynthetic' | ...
  x1: number; y1: number;  // tail point (canvas px)
  x2: number; y2: number;  // head point (canvas px)
  cpx: number; cpy: number; // Bezier control point (for curved arrows)
  label?: string;           // reaction condition text, e.g. "H2O, 80C"
}

interface CanvasState {
  atoms: Atom[];
  bonds: Bond[];
  arrows: Arrow[];
}
\end{lstlisting}

\lstinline|CanvasState| is the single source of truth. Everything else --- SMILES, the
2D SVG preview, the 3D viewer, CDXML output --- is derived from this. Think of it like
a Julia \lstinline|struct| that you pass to all downstream functions.

The spectrum prediction panels define their own result types:

\begin{lstlisting}[language=TypeScript]
// A single NMR chemical shift prediction for one (set of equivalent) nucleus/nuclei
interface NmrPeak {
  shift: number;    // chemical shift in ppm
  atom_idx: number; // 0-based heavy-atom index in the RDKit molecule
  count: number;    // number of equivalent nuclei (e.g. 3 for a methyl 1H)
}

// Full NMR prediction result returned by the sidecar's predict_nmr command
interface NmrSpectrum {
  peaks_1h:  NmrPeak[];   // proton shifts (H atoms grouped by attached heavy atom)
  peaks_13c: NmrPeak[];   // carbon shifts (one entry per heavy C atom)
}
\end{lstlisting}

\textbf{A note on \lstinline|electrons|}: stores lone pair electrons for display (shown
as dots). Zero means no dots, 2 means one pair, 4 means two pairs. Purely cosmetic ---
RDKit ignores it when generating SMILES.

\textbf{On \lstinline|element| vs \lstinline|label|}: \lstinline|element| is what the
atom \textit{is} chemically. For shorthand groups like \lstinline|OMe| or \lstinline|Ph|,
\lstinline|element| stores the shorthand string and \texttt{graph.ts} has a lookup table
(\lstinline|SHORTHAND_DATA|) that maps each shorthand to its heavy-atom lead element
for the V2000 Molblock.

\subsection{7b. Global State (src/store/index.ts)}

The Zustand store holds all application state in one place, divided into four logical
groups.

\textbf{Canvas state} --- the chemical graph and its history:
\begin{lstlisting}[language=TypeScript]
atoms: Atom[];
bonds: Bond[];
arrows: Arrow[];
history: CanvasState[];    // undo/redo snapshots
historyIndex: number;      // current position in history
selectedAtomIds: Set<string>;
selectedArrowIds: Set<string>;
lastInteractedAtomId: string | null;  // used for viewer focus (see 7e)
canvasSize: { width: number; height: number };
\end{lstlisting}

\textbf{Editor state} --- current tool settings:
\begin{lstlisting}[language=TypeScript]
tool: Tool;           // 'bond' | 'atom' | 'select' | 'eraser' | ...
element: string;      // currently selected atom element ("C", "N", etc.)
bondOrder: number;    // 1, 2, or 3
stereo: number;       // 0=plain, 1=wedge, 6=hash
arrowType: ArrowType;
currentCharge: number;
lonePairMode: boolean;    // charge tool placing lone pairs, not charges
selectedFragment: string; // SMILES of current ring template
showHydrogens: boolean;
\end{lstlisting}

\textbf{Viewer state} --- 2D/3D preview panel:
\begin{lstlisting}[language=TypeScript]
viewerSmiles: string;    // SMILES currently shown in the viewer
editableSmiles: string;  // SMILES in the text input bar (can differ while typing)
isInputFocused: boolean; // whether the SMILES text input has focus (blocks hotkeys)
previewMode: '2D' | '3D';
showViewer: boolean;
viewerMode: ViewerMode;  // 'pinned' | 'floating' | 'split'
viewerPos: { x: number; y: number };
viewerSize: { width: number; height: number };
splitWidth: number;
nmrSpectrum: NmrSpectrum | null;  // most recent NMR prediction (null until first prediction)
\end{lstlisting}

\textbf{Undo/redo} works by maintaining a snapshot list.
\lstinline|pushToHistory| appends the new \lstinline|CanvasState| to
\lstinline|history| (truncating any redo branch) and atomically updates the live
atoms/bonds/arrows. \lstinline|undo| decrements \lstinline|historyIndex| and restores
that snapshot:

\begin{lstlisting}[language=TypeScript]
pushToHistory: (newState) => set((s) => {
  const hist = s.history.slice(0, s.historyIndex + 1); // discard redo branches
  hist.push(newState);
  if (hist.length > 50) hist.shift(); // keep at most 50 levels
  return {
    history: hist,
    historyIndex: hist.length - 1,
    atoms: newState.atoms,
    bonds: newState.bonds,
    arrows: newState.arrows,
  };
}),

undo: () => set((s) => {
  if (s.historyIndex > 0) {
    const prev = s.history[s.historyIndex - 1];
    return { historyIndex: s.historyIndex - 1, ...prev };
  }
  if (s.historyIndex === 0)
    return { historyIndex: -1, atoms: [], bonds: [], arrows: [] };
  return {};
}),
\end{lstlisting}

There are two distinct ways to update the canvas state:
\begin{itemize}[leftmargin=1.5em]
  \item \lstinline|pushToHistory(state)| --- commits a change, adds it to the undo
        stack. Use for user-driven actions like drawing a bond or moving atoms.
  \item \lstinline|setCurrentCanvasState(state)| --- updates the live state
        \textit{without} pushing to history. Use for in-progress mouse-drag previews
        that shouldn't create intermediate undo entries.
\end{itemize}

\subsection{7c. App.tsx: Root Component}

\lstinline|App.tsx| is the outermost React component. It owns layout, menus, keyboard
shortcuts, file operations, the viewer layout system, and the window controls. It does
not do chemistry --- it delegates all of that to \lstinline|ChemCanvas|.

\subsubsection{The Theme System}

There are no CSS theme files, no CSS variables, and no theme library. The theme is a
plain JavaScript object:

\begin{lstlisting}[language=TypeScript]
const theme = {
  bg:      isDarkMode ? '#1e1e1e' : '#ebebeb',
  sidebar: isDarkMode ? '#252526' : '#f3f3f3',
  header:  isDarkMode ? '#333333' : '#ffffff',
  text:    isDarkMode ? '#cccccc' : '#444444',
  border:  isDarkMode ? '#444444' : '#dcdcdc',
  canvas:  isDarkMode ? '#1e1e1e' : '#f5f5f5',
};
\end{lstlisting}

This object is then passed directly into inline \lstinline|style| props throughout the
JSX:
\begin{lstlisting}[language=TypeScript]
<div style={{ background: theme.bg, color: theme.text,
              border: `1px solid ${theme.border}` }}>
\end{lstlisting}

The advantage: no CSS specificity to fight, no class name collisions, and the
relationship between \lstinline|isDarkMode| and every colour is explicit and local.

\subsubsection{The Frameless Window}

\lstinline|tauri.conf.json| sets \lstinline|"decorations": false|, removing the OS
title bar. The header div has \lstinline|data-tauri-drag-region|:
\begin{lstlisting}[language=TypeScript]
<div data-tauri-drag-region style={{ background: theme.header, height: '36px' }}>
\end{lstlisting}

This attribute is a Tauri hook. Any element with it can be used to drag the window ---
Tauri intercepts mouse events on such elements and routes them to the OS window manager.

The window control buttons call into Tauri's window API:
\begin{lstlisting}[language=TypeScript]
import { getCurrentWindow } from '@tauri-apps/api/window';
const appWindow = getCurrentWindow();

<button onClick={() => appWindow.minimize()}>        // minimize
<button onClick={() => appWindow.toggleMaximize()}>  // maximize
<button onClick={() => appWindow.close()}>           // close
\end{lstlisting}

\subsubsection{Keyboard Shortcuts}

\lstinline|App.tsx| registers global keyboard shortcuts via a \lstinline|useEffect|
with a \lstinline|keydown| listener on \lstinline|window|:

\begin{lstlisting}[language=TypeScript]
useEffect(() => {
  const handler = (e: KeyboardEvent) => {
    // Guard: if the SMILES text input is focused, let keystrokes go to it
    if (useStore.getState().isInputFocused) return;

    if (e.ctrlKey || e.metaKey) {
      if (e.key === 's') { e.preventDefault(); handleSaveNative(); }  // Ctrl+S
      if (e.key === 'o') { e.preventDefault(); handleOpen(); }        // Ctrl+O
      if (e.key === 'z') { e.preventDefault(); undo(); }              // Ctrl+Z
      if (e.key === 'y') { e.preventDefault(); redo(); }              // Ctrl+Y
      if (e.key === 'a') { e.preventDefault(); selectAll(); }         // Ctrl+A
      if (e.key === 'k') { e.preventDefault(); canvasRef.current?.cleanUp(); }
    }
    if (e.key === 'Escape') { deselectAll(); closeAllMenus(); }
  };
  window.addEventListener('keydown', handler);
  return () => window.removeEventListener('keydown', handler);
}, [handleSaveNative, handleOpen]);
\end{lstlisting}

The \lstinline|isInputFocused| guard is important: without it, pressing \texttt{C} to
switch to carbon would also type ``C'' into the SMILES text box when it's focused.

Additional hotkeys live in \lstinline|ChemCanvas.tsx|:

\begin{tabularx}{\linewidth}{@{} l X @{}}
\toprule
\textbf{Key(s)} & \textbf{Action} \\
\midrule
\texttt{c n o f s p b i l h} & Switch hovered atom to C, N, O, F, S, P, Br, I, Cl, H \\
\texttt{m e a d} & Switch hovered atom to Me, Et, Ph, D (deuterium) \\
\texttt{1 2 3} & Set bond order to single/double/triple \\
\texttt{w} & Wedge bond (bold) \\
\texttt{4 5 6 7 8} & Place cyclobutane/cyclopentane/cyclohexane/etc.\ ring \\
\texttt{a} & Place benzene ring \\
\texttt{Delete / Backspace} & Remove selected atoms and their bonds \\
\bottomrule
\end{tabularx}

\subsubsection{The Viewer Layout Modes}

The 2D/3D viewer can appear in three modes controlled by \lstinline|viewerMode|:

\begin{itemize}[leftmargin=1.5em]
  \item \textbf{\lstinline|split|} --- The viewer occupies a right panel; a draggable
        divider separates it from the canvas. The canvas width shrinks accordingly.
  \item \textbf{\lstinline|pinned|} --- The viewer is an overlay panel docked to the
        right edge of the canvas with a resize handle.
  \item \textbf{\lstinline|floating|} --- The viewer is a free-floating overlay that
        can be dragged anywhere on the canvas.
\end{itemize}

\subsubsection{The SMILES Bar}

There are two distinct SMILES strings in the store:
\begin{itemize}[leftmargin=1.5em]
  \item \lstinline|viewerSmiles| --- the SMILES of the currently focused fragment,
        drives the 2D/3D viewer. Updated automatically as atoms are drawn.
  \item \lstinline|editableSmiles| --- the SMILES in the text input. Normally mirrors
        \lstinline|viewerSmiles|, but can diverge while the user types. Pressing Enter
        calls \lstinline|loadFromSmiles()| on the canvas, which parses the SMILES via
        RDKit and loads the result onto the canvas.
\end{itemize}

\subsubsection{Fragment Ring Templates}

The sidebar has a RING section with buttons for common ring systems. Each button calls:
\begin{lstlisting}[language=TypeScript]
setTool('fragment');
setSelectedFragment('c1ccccc1');  // benzene, or C1CCCCC1, etc.
\end{lstlisting}

When a fragment is placed, \lstinline|placeFragmentAt()| in \lstinline|ChemCanvas.tsx|
handles the geometry:
\begin{enumerate}[leftmargin=1.5em]
  \item Ask RDKit to generate a fresh 2D embedding of the SMILES;
        parse the resulting atom coordinates from the molblock.
  \item If clicking on an \textbf{existing bond}: align the fragment so that
        edge[0]--edge[1] maps onto that bond's vector; rotate to minimise overlap.
  \item If clicking on an \textbf{existing atom}: align fragment[0] to that atom;
        choose the rotation that points into the largest angular gap in existing bonds.
  \item If clicking on \textbf{empty space}: place the fragment centred at the click.
  \item Merge overlapping atoms (the atom you clicked becomes the ring junction),
        add new atoms and bonds, push to history.
\end{enumerate}

\subsection{7d. The Drawing Engine (ChemCanvas.tsx)}

\lstinline|ChemCanvas.tsx| owns the canvas rendering, hit testing, zoom/pan,
keyboard handling, and the SMILES update pipeline.

\subsubsection{react-konva}

Konva is a Canvas 2D drawing library. react-konva provides React-style JSX bindings.
Instead of calling the Canvas 2D API imperatively:

\begin{lstlisting}[language=TypeScript]
// Plain Canvas 2D API (without react-konva)
ctx.beginPath();
ctx.arc(atom.x, atom.y, 8, 0, Math.PI * 2);
ctx.fillStyle = "white";
ctx.fill();
ctx.stroke();
\end{lstlisting}

You write declarative JSX:
\begin{lstlisting}[language=TypeScript]
<Circle x={atom.x} y={atom.y} radius={8} fill="white" stroke="black" />
\end{lstlisting}

The Konva rendering tree:
\begin{lstlisting}[language=TypeScript]
<Stage width={width} height={height}>    // the canvas element
  <Layer>                                // a drawing layer
    <Line points={[x1,y1, x2,y2]} stroke="black" strokeWidth={2} />
    <Circle x={50} y={50} radius={8} fill="white" />
    <Text x={50} y={50} text="C" />
    <Shape sceneFunc={(ctx, shape) => { /* raw Canvas 2D */ }} />
  </Layer>
</Stage>
\end{lstlisting}

\lstinline|Shape| with \lstinline|sceneFunc| is the escape hatch for drawing wedge bonds
and arrows that require complex paths.

\subsubsection{Zoom and Pan}

The entire canvas contents are wrapped in a Konva transform:
\begin{lstlisting}[language=TypeScript]
const [stageScale, setStageScale] = useState(1);
const [stagePos, setStagePos] = useState({ x: 0, y: 0 });
\end{lstlisting}

All atom coordinates are passed as-is to Konva; the stage applies
\lstinline|scaleX={stageScale} scaleY={stageScale} x={stagePos.x} y={stagePos.y}|
globally. When converting mouse pixel coordinates to logical coordinates:
\begin{lstlisting}[language=TypeScript]
const getRelativePointerPosition = () => {
  const pointer = stage.getPointerPosition();   // pixel position on screen
  return {
    x: (pointer.x - stagePos.x) / stageScale,  // logical x
    y: (pointer.y - stagePos.y) / stageScale,  // logical y
  };
};
\end{lstlisting}

The hit radius for atoms scales inversely with zoom:
\begin{lstlisting}[language=TypeScript]
const findAtomAt = (x: number, y: number) => {
  const r = 15 / stageScale;   // constant pixel radius regardless of zoom
  return atoms.find(a => Math.sqrt((a.x-x)**2 + (a.y-y)**2) < r) ?? null;
};
\end{lstlisting}

\subsubsection{Hit Testing}

Three hit-test functions:

\textbf{\lstinline|findAtomAt(x, y)|} --- Euclidean distance from cursor to atom
centre, within radius \lstinline|15/scale|.

\textbf{\lstinline|findBondAt(x, y)|} --- Point-to-line-segment distance. For each
bond, project the cursor onto the segment clamped to $[0,1]$, compute distance, return
the nearest bond within \lstinline|10/scale|.

\textbf{\lstinline|findArrowAt(x, y)|} --- For straight arrows: same point-to-segment
as bonds. For curved arrows: sample 20 points along the Bezier curve
($t = 0, 0.05, \ldots, 1.0$) and check if any sample is within \lstinline|12/scale|
of the cursor.

\subsubsection{useBondTool in Detail}

Bond drawing has three phases:

\textbf{Mouse down:} If you clicked on an existing atom, start dragging from it. If on
an existing bond, cycle or set its order/stereo. If on empty space, create a new carbon
atom at that position and start dragging from it.

\textbf{Mouse move:} Compute the current bond preview. The angle is snapped to
$15^\circ$ increments ($\pi/12$ radians):
\begin{lstlisting}[language=TypeScript]
const dx = pos.x - start.x, dy = pos.y - start.y;
const angle = Math.round(Math.atan2(dy, dx) / (Math.PI / 12)) * (Math.PI / 12);
\end{lstlisting}

\textbf{Mouse up:} If the drag was shorter than 5 pixels (essentially a click), auto-place
a bond in the chemically natural direction --- find the largest angular gap between
existing bonds and bisect it:
\begin{lstlisting}[language=TypeScript]
const gaps = [];
for (let i = 0; i < angles.length; i++) {
  let gap = angles[(i+1) % angles.length] - angles[i];
  if (gap <= 0) gap += 2 * Math.PI;
  gaps.push({ gap, bisector: angles[i] + gap / 2 });
}
const best = gaps.sort((a, b) => b.gap - a.gap)[0];
finalAngle = best.bisector;
\end{lstlisting}

\subsubsection{useSelectTool in Detail}

The select tool handles four distinct interactions:
\begin{enumerate}[leftmargin=1.5em]
  \item \textbf{Click on atom} --- select it (or shift-click to add/remove from selection)
  \item \textbf{Click on bond} --- if both endpoints are already selected, cycle bond
        order; otherwise select both endpoints
  \item \textbf{Drag selected atoms} --- translate all selected atoms by the mouse
        delta, using \lstinline|setCurrentCanvasState| (no history entry until mouse-up)
  \item \textbf{Rubber-band selection} --- click on empty space, drag to draw a box,
        select all atoms within the box
\end{enumerate}

There is also a \textbf{rotation handle}: when atoms are selected, a small circle appears
60\,px above the centroid. Dragging it rotates the selection around its centroid. With
Shift held, rotation snaps to $15^\circ$ increments. During drag,
\lstinline|setCurrentCanvasState| is called on every mouse move; on mouse-up,
\lstinline|pushToHistory| commits the final state as a single undo entry.

\subsubsection{useAtomTool in Detail}

The atom tool behaves similarly to the bond tool, but clicking on an existing atom
(without dragging) opens an \textbf{inline text editor}:

\begin{lstlisting}[language=TypeScript]
if (dist < 5 && !preview.targetId) {
  if (activeAtomWasExisting.current) {
    setEditingAtomId(start.id);   // triggers a text input in ChemCanvas
    setEditingValue(start.element);
  }
}
\end{lstlisting}

\lstinline|ChemCanvas| renders a positioned \lstinline|<input>| element over the atom
when \lstinline|editingAtomId| is set. Pressing Enter commits the new label and pushes
to history. This lets you type arbitrary shorthand labels directly onto atoms.

\subsection{7e. SMILES Generation Pipeline}

\begin{lstlisting}[language={}]
canvas atoms/bonds
      |
      v  hashCanvasState(atoms, bonds)  -- cheap string hash
      |  if hash unchanged: skip (memoization)
      |  (200ms debounce: only runs 200ms after the last change)
      |
      v  graphToMolblock(atoms, bonds, width, height)  [src/lib/graph.ts]
      |
      v  V2000 Molblock string
      |
      +---> rdkit.get_mol(molblock) -> mol.get_smiles() -> editableSmiles (Structure Input bar)
      |
      v  getConnectedComponents(atoms, bonds)
      |  identify which fragment to show in the viewer
      |  (priority: selected atom's fragment > last interacted > first component)
      |
      v  graphToMolblock(fragment atoms/bonds) -> rdkit.get_mol() -> mol.get_smiles()
      |
      v  viewerSmiles -> Molecule.tsx (2D SVG) + Molecule3DThreePanel.tsx (3D)
\end{lstlisting}

The SMILES update is triggered by a \lstinline|useEffect| watching \lstinline|atoms|,
\lstinline|bonds|, \lstinline|rdkit|, \lstinline|width|, and \lstinline|height|. It is
\textbf{debounced by 200\,ms}: a \lstinline|setTimeout| is set on each change, and
the previous timer is cancelled via the \lstinline|useEffect| cleanup return. SMILES is
not recomputed while the user is actively dragging.

\subsubsection{hashCanvasState}

\begin{lstlisting}[language=TypeScript]
export function hashCanvasState(atoms: Atom[], bonds: Bond[]): string {
  let h = `${atoms.length}:${bonds.length}`;
  for (const a of atoms)
    h += `|${a.id}${a.x.toFixed(1)}${a.y.toFixed(1)}${a.element}${a.charge ?? 0}`;
  for (const b of bonds)
    h += `|${b.from}${b.to}${b.order}${b.stereo ?? 0}`;
  return h;
}
\end{lstlisting}

This is a cheap string concatenation hash. The \lstinline|.toFixed(1)| means positions
that differ by less than 0.1\,px are treated as identical (minor floating-point drift
won't trigger SMILES recomputation).

\subsubsection{getConnectedComponents and Viewer Focus}

When there are multiple disconnected structures on the canvas (e.g.\ a reaction scheme),
the \textbf{viewer} should show just one fragment --- whichever the user is working on.
\lstinline|getConnectedComponents| finds the connected components via BFS:

\begin{lstlisting}[language=TypeScript]
export function getConnectedComponents(atoms: Atom[], bonds: Bond[]): Set<string>[] {
  const adj = buildAdjacencyMap(atoms, bonds);
  const visited = new Set<string>();
  const components: Set<string>[] = [];
  for (const a of atoms) {
    if (visited.has(a.id)) continue;
    const comp = new Set<string>();
    const q = [a.id];
    while (q.length) {
      const curr = q.shift()!;
      comp.add(curr);
      for (const n of adj.get(curr) ?? []) {
        if (!visited.has(n)) { visited.add(n); q.push(n); }
      }
    }
    components.push(comp);
  }
  return components;
}
\end{lstlisting}

Viewer focus logic: (1) if any atom is selected, show its fragment; (2) else if there's
a \lstinline|lastInteractedAtomId|, show its fragment; (3) else show the first component.

\subsubsection{graphToMolblock and the V2000 Format}

The key transformation is coordinate scaling:

\begin{lstlisting}[language=TypeScript]
// Canvas pixels -> Angstroms, centred at the canvas midpoint
const x_ang = (atom.x - canvasWidth  / 2) / MOLBLOCK_SCALE;  // MOLBLOCK_SCALE = 30
const y_ang = -(atom.y - canvasHeight / 2) / MOLBLOCK_SCALE;  // Y flipped: canvas Y grows down
\end{lstlisting}

\texttt{MOLBLOCK\_SCALE = 30} means 30 canvas pixels $= 1$\,\AA. The V2000 format is
fixed-width text:

\begin{lstlisting}[language={}]
\n  ChemEditor\n\n
  3  2  0  0  0  0  0  0  0  0999 V2000
    0.0000    0.0000    0.0000 C   0  0  0  0  0  0  0  0  0  0  0  0
    1.5000    0.0000    0.0000 O   0  0  0  0  0  0  0  0  0  0  0  0
  1  2  1  0  0  0  0   <- bond: atom1, atom2, order, stereo
M  CHG  1   2  -1       <- charge block
M  END
\end{lstlisting}

\textbf{V2000 stereo codes} (MDL-format integers):
\begin{itemize}[leftmargin=1.5em]
  \item \lstinline|stereo = 0| --- no stereo (plain line)
  \item \lstinline|stereo = 1| --- wedge bond (bold triangle, towards viewer)
  \item \lstinline|stereo = 6| --- hash bond (dashed lines, away from viewer)
\end{itemize}

\textbf{Shorthand handling}: \lstinline|SHORTHAND_DATA| maps labels like
\lstinline|"OMe"|, \lstinline|"Ph"|, \lstinline|"tBu"| to their heavy-atom lead
element. When writing the Molblock:
\begin{lstlisting}[language=TypeScript]
const elem = SHORTHAND_DATA[atom.element]?.lead ?? atom.element;
// "OMe" -> "O", "Ph" -> "C", "tBu" -> "C", "N" -> "N"
\end{lstlisting}

\subsubsection{molblockToState --- The Inverse Function}

\lstinline|molblockToState| is the reverse of \lstinline|graphToMolblock|. It is used
when loading files and when the user types a SMILES into the bar:

\begin{lstlisting}[language=TypeScript]
export function molblockToState(
  mb: string,
  canvasWidth: number,
  canvasHeight: number,
  existingArrows: Arrow[] = []
): CanvasState {
  const ls = mb.split('\n');
  const cL = ls[3];
  const nA = parseInt(cL.substring(0, 3).trim());
  const nB = parseInt(cL.substring(3, 6).trim());

  const atoms: Atom[] = [];
  for (let i = 0; i < nA; i++) {
    const l = ls[i + 4];
    atoms.push({
      id: crypto.randomUUID(),  // fresh UUIDs for all atoms
      x: canvasWidth  / 2 + parseFloat(l.substring(0,  10)) * MOLBLOCK_SCALE,
      y: canvasHeight / 2 - parseFloat(l.substring(10, 20)) * MOLBLOCK_SCALE,
      element: l.substring(31, 34).trim(),
    });
  }
  // ... parse bonds, charge block ...
  return { atoms, bonds, arrows: existingArrows };
}
\end{lstlisting}

Note that \lstinline|existingArrows| are passed through unchanged --- MOL/SDF files
have no concept of reaction arrows, so existing arrows are preserved.

\subsection{7f. 2D Preview (Molecule.tsx)}

\lstinline|Molecule.tsx| takes a SMILES string and renders it as an image.
The RDKit WASM API generates SVG:

\begin{lstlisting}[language=TypeScript]
const svg = useMemo(() => {
  if (!rdkit || !smiles) return null;
  const mol = rdkit.get_mol(smiles);
  if (!mol) return "Invalid SMILES";

  const params = {
    width: Math.round(width),
    height: Math.round(height),
    backgroundColour: [1, 1, 1, 0],   // transparent (RGBA, 0-1 range)
    bondLineWidth: 2.0,
    symbolColour: isDarkMode ? [1, 1, 1, 1] : [0, 0, 0, 1],
    ...(isDarkMode && { atomColourPalette: { 6: [0.75, 0.75, 0.75] } }),
    // ^ In dark mode, override C (element 6) to grey so bonds are visible
  };

  let svgString = mol.get_svg_with_highlights(JSON.stringify(params));
  mol.delete();   // must free WASM heap memory manually
  return svgString;
}, [rdkit, smiles, width, height, isDarkMode]);
\end{lstlisting}

The SVG string is then turned into an image using the \textbf{Blob URL} pattern:

\begin{lstlisting}[language=TypeScript]
useEffect(() => {
  if (!svg || !svg.startsWith('<')) { setImgUrl(null); return; }
  const blob = new Blob([svg], { type: 'image/svg+xml' });
  const url = URL.createObjectURL(blob);
  setImgUrl(url);
  return () => URL.revokeObjectURL(url);  // cleanup when svg changes
}, [svg]);

// In JSX:
<img src={imgUrl} width={width} height={height} />
\end{lstlisting}

\textbf{Why Blob URL instead of \lstinline|dangerouslySetInnerHTML|?} The Blob URL
approach creates a temporary \lstinline|blob:| URL that the browser loads as an image.
This handles SVG rendering more consistently, and \lstinline|URL.revokeObjectURL| cleans
up memory cleanly when SMILES changes. Without \lstinline|revokeObjectURL|, every SMILES
change would leave an orphaned Blob in browser memory.

\subsubsection{The useRDKit Singleton}

\begin{lstlisting}[language=TypeScript]
let rdkitPromise: Promise<any> | null = null;   // module-level variable

export function useRDKit() {
  const [rdkit, setRdkit] = useState<any>(null);
  const [loading, setLoading] = useState(true);

  useEffect(() => {
    if (!rdkitPromise) {
      // First call: start loading WASM. Store the Promise at module level.
      rdkitPromise = window.initRDKitModule({
        locateFile: () => '/rdkit/RDKit_minimal.wasm',
      });
    }
    // All subsequent calls: attach to the already-started promise.
    rdkitPromise.then((instance) => {
      setRdkit(instance);
      setLoading(false);
    });
  }, []);

  return { rdkit, loading };
}
\end{lstlisting}

The module-level \lstinline|rdkitPromise| variable is what makes this a singleton. No
matter how many components call \lstinline|useRDKit()| ---
\lstinline|ChemCanvas|, \lstinline|Molecule|, \lstinline|App| --- the WASM is only loaded
once. Without this, each component would start its own $\sim$1\,second WASM
initialisation.

\subsection{7g. 3D Viewer (Molecule3DThreePanel.tsx + Molecule3DThree.tsx)}

The 3D viewer is now split into two layers:
\begin{itemize}[leftmargin=1.5em]
  \item \texttt{Molecule3DThreePanel.tsx} handles viewer orchestration: conformer
        loading, viewer controls, the conformer strip, and the \texttt{forwardRef}
        bridge used by \texttt{App.tsx}.
  \item \texttt{Molecule3DThree.tsx} is the low-level \texttt{three.js} renderer:
        scene setup, lighting, camera controls, mesh generation, highlighting,
        measurement interaction, and screenshot capture.
\end{itemize}

\subsubsection{The Full Data Flow}

\begin{lstlisting}[language={}]
viewerSmiles (SMILES string)
      |
      v  useMolecule3DData({ smiles, forceField, maxConformers, ... })
      |
      v  fetchConformers(smiles, forceField, limit)
      |    Check module-level cache: Map<"smiles|ff|max", Viewer3DMolecule[]>
      |    Cache miss -> invoke('generate_conformers', { smiles, ... })
      |               -> Rust IPC -> Python sidecar subprocess
      |
      v  Python returns JSON:
      |  {
      |    conformers: [
      |      { atoms: [{x,y,z,element}, ...], energy: -45.2, delta_e: 0.0 },
      |      ...
      |    ],
      |    bonds: [{a1: 0, a2: 1, order: 1}, ...]  // shared across conformers
      |  }
      |
      v  alignViewerMoleculeToReference(previous, next)
      |    Keep conformer switches visually stable when topology matches
      |
      v  Molecule3DThree.tsx
      |    build atom/bond meshes
      |    fit the camera to the visual bounds
      |    update the directional light rig relative to the camera
      |
      v  WebGL rendered in the containerRef <div>
\end{lstlisting}

\subsubsection{Conformer Caching}

The cache still lives at module level (outside the component function), so it persists
across renders and even if the viewer is closed and reopened. The current
implementation also keeps an in-flight promise cache so repeated requests for the same
structure do not start duplicate sidecar work:

\begin{lstlisting}[language=TypeScript]
const conformerCache = new Map<string, Viewer3DMolecule[]>();
const inFlightConformerCache = new Map<string, Promise<Viewer3DMolecule[]>>();

async function fetchConformers(smiles, forceField, maxConformers) {
  const key = `${smiles}|${forceField}|${maxConformers}`;
  const cached = conformerCache.get(key);
  if (cached) return cached;
  if (inFlightConformerCache.has(key)) return inFlightConformerCache.get(key)!;

  const request = invoke('generate_conformers', { smiles, forceField, maxConformers });
  inFlightConformerCache.set(key, request);
  // ... process result ...
}
\end{lstlisting}

\subsubsection{Orientation Stability}

The current viewer keeps orientation stable in two ways:
\begin{itemize}[leftmargin=1.5em]
  \item When the same topology is shown with a different conformer, coordinates are
        pre-aligned to the previous conformer before rendering.
  \item When the topology changes, the camera is re-fit to the new molecule's visual
        bounds.
\end{itemize}

\begin{lstlisting}[language=TypeScript]
if (viewer3dHasMatchingTopology(previousMolecule, molecule)) {
  molecule = alignViewerMoleculeToReference(previousMolecule, molecule);
}
\end{lstlisting}

\subsubsection{Three.js Scene Setup}

The renderer is initialised imperatively:
\begin{lstlisting}[language=TypeScript]
const renderer = new THREE.WebGLRenderer({
  antialias: true,
  alpha: false,
  preserveDrawingBuffer: true,
});
const scene = new THREE.Scene();
const camera = new THREE.PerspectiveCamera(fov, aspect, 0.01, 1000);
const controls = new TrackballControls(camera, renderer.domElement);
\end{lstlisting}

The scene also owns a camera-relative light rig (ambient, hemisphere, key, fill, rim)
so the model stays readable as the user rotates it.

\subsubsection{captureViewer via forwardRef}

\lstinline|Molecule3DThree| uses \lstinline|forwardRef| to expose
\lstinline|captureViewer()| to its parent:
\begin{lstlisting}[language=TypeScript]
useImperativeHandle(ref, () => ({
  captureViewer: () => {
    const originalPixelRatio = renderer.getPixelRatio();
    const exportPixelRatio = Math.min(4, Math.max(originalPixelRatio, 3));
    renderer.setPixelRatio(exportPixelRatio);
    renderer.render(scene, camera);
    const image = renderer.domElement.toDataURL('image/png');
    renderer.setPixelRatio(originalPixelRatio);
    return image;
  },
}));
\end{lstlisting}

The temporary pixel-ratio bump makes exported PNGs much sharper than the live viewport.

\subsubsection{Measurement Tools}

Measurements are implemented directly with Three.js raycasting and overlay groups:
\begin{lstlisting}[language=TypeScript]
raycaster.setFromCamera(pointer, camera);
const atomHit = raycaster.intersectObjects(atomMeshes, false)[0];
// Collect atoms for distance / angle / dihedral,
// then draw highlight geometry into a dedicated overlay group.
\end{lstlisting}

The measurement visuals are not DOM overlays; they are actual Three.js objects that
share the molecule's coordinate frame.

\subsubsection{Resize Handling}

The viewer still needs explicit resize handling because the WebGL renderer and camera
aspect must track the panel size:

\begin{lstlisting}[language=TypeScript]
useEffect(() => {
  renderer.setSize(width, height);
  camera.aspect = width / height;
  camera.updateProjectionMatrix();
  controls.handleResize?.();
}, [width, height]);
\end{lstlisting}

This resize/update pass also reapplies the viewer background and re-fits the camera
after representation or panel-size changes.

\subsection{7h. The Python Sidecar}

\subsubsection{What a Sidecar Is}

A ``sidecar'' in Tauri terminology is an external binary that Tauri manages as a
subprocess. It is bundled inside the final app package and Tauri handles finding and
spawning it. This is the mechanism for running code in languages other than
Rust/JavaScript from within a Tauri app.

\subsubsection{The Naming Convention --- Platform Triples}

Tauri requires the sidecar binary to be named with the full Rust target triple appended:

\begin{lstlisting}[language={}]
chem-engine-x86_64-unknown-linux-gnu     (Linux x86_64)
chem-engine-aarch64-unknown-linux-gnu    (Linux ARM64)
chem-engine-x86_64-apple-darwin         (macOS Intel)
chem-engine-aarch64-apple-darwin        (macOS Apple Silicon)
chem-engine-x86_64-pc-windows-msvc.exe  (Windows)
\end{lstlisting}

Find your machine's triple:
\begin{lstlisting}[language=bash]
rustc -vV | grep host
# host: x86_64-unknown-linux-gnu
\end{lstlisting}

If the file is named incorrectly, \lstinline|app.shell().sidecar("chem-engine")| returns
\lstinline|Err(...)| and the Rust code silently falls through to the development fallback
(\lstinline|uv run python|). There's no crash, just a missing binary.

\subsubsection{build-sidecar.js}

\lstinline|build-sidecar.js| automates the full sidecar build pipeline in seven steps:

\begin{enumerate}[leftmargin=1.5em]
  \item \textbf{Detect target triple} from \lstinline|rustc -vV| (e.g.
        \lstinline|x86_64-unknown-linux-gnu|).
  \item \textbf{Ensure \texttt{src-tauri/bin/} exists}.
  \item \textbf{Download xtb} $\sim$60\,MB tarball from GitHub Releases (Linux x86\_64
        only; other platforms use the system xtb on PATH).
  \item \textbf{Build the HOSE NMR database} (\lstinline|nmr_hose_db.json|) from
        \lstinline|development/nmr_model/data/nmrshiftdb2withsignals.sd| when the source
        data is newer than the bundled JSON.
  \item \textbf{Incremental check}: skip PyInstaller if the sidecar binary is newer
        than \lstinline|chem-engine.py|, the xtb binary, and the bundled data files.
  \item \textbf{Run PyInstaller} via \lstinline|uv run --group build pyinstaller|, passing
        \lstinline|--add-binary| / \lstinline|--add-data| flags for xtb,
        \lstinline|nmr_refs.json|, and \lstinline|nmr_hose_db.json|.
  \item \textbf{Rename} \lstinline|dist/chem-engine| to
        \lstinline|bin/chem-engine-<triple>| (Tauri's required naming).
\end{enumerate}

PyInstaller's \lstinline|--onefile| flag bundles the Python interpreter, all
\lstinline|rdkit| / \lstinline|numpy| dependencies,
\lstinline|chem-engine.py|, the xtb binary, and the bundled NMR data files into a single
self-extracting binary ($\sim$130--150\,MB). On first run, PyInstaller extracts its
contents to a temporary directory (\lstinline|sys._MEIPASS|) --- which is why the first
invocation of the sidecar is slightly slower.

\subsubsection{The stdin/stdout JSON Protocol}

The sidecar communicates via stdin (input to the Python process) and stdout (output from
it). Trivially testable from the terminal:

\begin{lstlisting}[language=bash]
# Test conformer generation
uv run python src-tauri/bin/chem-engine.py conformer "CCO" --ff MMFF94s --max 3

# Test file parsing (molecule content goes on stdin)
cat molecule.mol | uv run python src-tauri/bin/chem-engine.py parse mol

# Test RXN export
cat molecule.mol | uv run python src-tauri/bin/chem-engine.py export rxn --arrows '[]'
\end{lstlisting}

\subsubsection{Conformer Generation Pipeline}

\begin{lstlisting}[language=Python]
def generate_conformers(smiles, pool_size=50, max_output=10,
                        energy_window=10.0, rmsd_threshold=0.5,
                        force_field='MMFF94s'):
    mol = Chem.MolFromSmiles(smiles)
    mol = Chem.AddHs(mol)           # explicit H for correct 3D geometry

    params = AllChem.ETKDGv3()
    params.randomSeed = 0xf00d      # reproducible results
    conf_ids = AllChem.EmbedMultipleConfs(mol, numConfs=pool_size, params=params)

    ff_results = AllChem.MMFFOptimizeMoleculeConfs(
        mol, mmffVariant='MMFF94s', maxIters=500)
    # returns list of (converged, energy) tuples

    # Filter: keep only converged conformers, sorted by energy
    data_pool = sorted(
        [{"conf_id": i, "energy": e}
         for i, (conv, e) in zip(conf_ids, ff_results) if conv == 0],
        key=lambda x: x["energy"]
    )

    # Deduplicate: keep conformers >0.5 A RMSD apart AND within 10 kcal/mol
    unique = [data_pool[0]]
    for item in data_pool[1:]:
        if item["energy"] - data_pool[0]["energy"] > energy_window: break
        is_unique = all(
            AllChem.GetBestRMS(mol, mol, item["conf_id"], kept["conf_id"])
            >= rmsd_threshold
            for kept in unique
        )
        if is_unique: unique.append(item)
        if len(unique) >= max_output: break
\end{lstlisting}

\subsubsection{NMR Prediction via HOSE Lookup (predict\_nmr)}

The NMR panel now uses a classical HOSE-code predictor rather than an ONNX neural model.
An offline build step converts NMRShiftDB2 into a compact JSON lookup database keyed by
canonical local atom environments. At runtime, the sidecar:

\begin{enumerate}[leftmargin=1.5em]
  \item Generates HOSE codes for each carbon atom and each proton-bearing heavy atom.
  \item Looks up the most specific matching sphere in \lstinline|nmr_hose_db.json|.
  \item Falls back through lower HOSE spheres as needed.
  \item Blends the HOSE median with deterministic chemistry corrections for aromatic,
        vinylic, benzylic, alpha-to-hetero, and carbonyl-adjacent environments.
  \item Applies fast multiplicity / J-coupling heuristics for \textsuperscript{1}H,
        including split methylene and simple alkene cases.
  \item Returns structured \lstinline|signals_1h| / \lstinline|signals_13c| payloads for
        the custom React plotter.
\end{enumerate}

\textbf{Runtime data files.}
PyInstaller sets \lstinline|sys._MEIPASS| to its unpacked temp directory.
The sidecar reads \lstinline|nmr_hose_db.json| and \lstinline|nmr_refs.json| from there in
production, or from \lstinline|src-tauri/bin/| in development.

\textbf{Offline database build.}
\lstinline|development/build_nmr_hose_db.py| reads
\lstinline|development/nmr_model/data/nmrshiftdb2withsignals.sd| and writes
\lstinline|src-tauri/bin/nmr_hose_db.json|. This step is triggered automatically by
\lstinline|pnpm sync:sidecar| and \lstinline|pnpm build:sidecar| when the source SDF is newer
than the bundled JSON.

\subsection{7i. ChemDraw Compatibility (src/utils/cdxml.ts)}

CDXML is ChemDraw's XML-based file format. The crucial coordinate convention:
ChemDraw uses \textbf{30 points per bond length unit}. Our canvas uses
\lstinline|MOLBLOCK_SCALE = 30| pixels per \AA. This is not a coincidence --- the
canvas scale was chosen to match ChemDraw. Canvas pixel coordinates map directly to
ChemDraw point coordinates without any rescaling:

\begin{lstlisting}[language=TypeScript]
// Canvas pixels === ChemDraw points (both use 30 units/Angstrom)
// Only an offset is needed to centre the canvas on the ChemDraw page:
const CDXML_CENTER_X = 306.0;  // ChemDraw default page centre
const CDXML_CENTER_Y = 396.0;

function atomToXML(atom: Atom, canvasWidth: number, canvasHeight: number): string {
  const cdx = CDXML_CENTER_X + (atom.x - canvasWidth  / 2);
  const cdy = CDXML_CENTER_Y + (atom.y - canvasHeight / 2);
  return `<n id="..." p="${cdx.toFixed(2)} ${cdy.toFixed(2)}" Element="..." ... />`;
}
\end{lstlisting}

% ─────────────────────────────────────────────────────────────────────────────
\section{Build System}
\label{sec:build}

\subsection{vite.config.ts}

\begin{lstlisting}[language=TypeScript]
export default defineConfig(async () => ({
  plugins: [react()],   // enables JSX compilation and fast refresh
  clearScreen: false,   // don't clear the terminal -- Tauri and Vite share it
  server: {
    port: 1420,         // Tauri expects the frontend at this port
    strictPort: true,   // fail if 1420 is in use rather than trying 1421, ...
    hmr: host ? { protocol: "ws", host, port: 1421 } : undefined,
    watch: { ignored: ["**/src-tauri/**"] },  // don't watch Rust files
  },
}));
\end{lstlisting}

\textbf{Why port 1420?} \texttt{tauri.conf.json} has
\lstinline|"devUrl": "http://localhost:1420"| --- this is the URL the Tauri WebView
loads in development mode. They must match.

\subsection{pnpm dev --- Development Mode, Step by Step}

\begin{lstlisting}[language={}]
1. "predev" script runs first (package.json):
   mkdir -p public/rdkit
   cp node_modules/@rdkit/rdkit/dist/RDKit_minimal.{js,wasm} public/rdkit/

   WHY: Vite only serves files from the project root. @rdkit/rdkit lives in
   node_modules/ which Vite does not serve directly. Copying to public/rdkit/
   makes these files available at /rdkit/RDKit_minimal.wasm in the browser.

2. "dev" script runs:
   vite

   Vite starts an HTTP dev server at localhost:1420.
   It processes TypeScript, JSX, CSS imports on demand.
   It establishes a WebSocket for Hot Module Replacement (HMR):
   when you save a .tsx file, Vite sends the changed module to the browser,
   and React updates only the changed component without a full page reload.

3. Separately, in another terminal:
   pnpm tauri:dev

   Cargo compiles src-tauri/src/ in debug mode (fast compile, slower runtime).
   Tauri opens a native window with WebView pointed at http://localhost:1420.
   Python sidecar: run_engine() falls back to "uv run python chem-engine.py"
   because the PyInstaller binary hasn't been built yet.
\end{lstlisting}

\subsection{pnpm build:web --- Frontend + Sidecar Assets}

\begin{lstlisting}[language={}]
1. pnpm build:sidecar  ->  node src-tauri/build-sidecar.js
   - Detects platform triple from `rustc -vV`
   - Runs: uv run pyinstaller --onefile bin/chem-engine.py --clean --noconfirm
     * PyInstaller finds all imports in chem-engine.py
     * Bundles Python interpreter + rdkit + all dependencies
     * Output: src-tauri/dist/chem-engine  (~120MB, self-extracting)
   - Renames: dist/chem-engine -> bin/chem-engine-x86_64-unknown-linux-gnu

2. tsc
   - TypeScript type checker only. Does NOT produce JavaScript output.
     (Vite handles transpilation; tsc here is just a type-safety gate.)
   - Fails the build if any type errors are found.

3. vite build
   - Bundles all TypeScript/JSX into optimised ES modules.
   - Applies tree-shaking (removes unused imports).
   - Minifies JavaScript and CSS.
   - Output: dist/ directory with index.html and hashed asset files.
\end{lstlisting}

\subsection{pnpm build / pnpm build:app / pnpm tauri:build --- Final Distributable}

\begin{lstlisting}[language={}]
pnpm build and pnpm build:app both resolve to pnpm tauri:build, which runs:

node scripts/tauri-build.mjs
   - On Linux, defaults to: pnpm tauri build --bundles deb,rpm
   - If CHEM_EDITOR_BUNDLE_TARGETS is set, forwards that bundle list directly
   - On macOS and Windows, runs the default Tauri bundle target set

cargo tauri build, which:

1. Runs the full Tauri build path using the current frontend dist and packaged sidecar.

2. Compiles src-tauri/src/ with rustc --release
   - Full optimisations, link-time optimisation (LTO)
   - Takes 2-5 minutes on first build (Cargo caches subsequent builds)
   - Output: src-tauri/target/release/chem-editor  (~5MB stripped binary)

3. Packages everything into distributable formats:
   src-tauri/target/release/bundle/deb/chem-editor_0.1.0_amd64.deb
   src-tauri/target/release/bundle/rpm/chem-editor-0.1.0-1.x86_64.rpm

   The bundle contains:
   - The Rust binary (chem-editor)
   - The built frontend (dist/ -- HTML, JS, CSS)
   - The PyInstaller sidecar (bin/chem-engine-x86_64-unknown-linux-gnu)
   - App icons
\end{lstlisting}

The final \texttt{.deb} installs ChemEditor as a standard desktop application. The
PyInstaller binary is self-contained --- users do not need Python or RDKit installed.

% ─────────────────────────────────────────────────────────────────────────────
\section{File I/O Flow}
\label{sec:fileio}

\subsection{Opening a File}

\begin{lstlisting}[language={}]
User clicks File -> Open
      |
      v  open({ multiple: false, filters: [
           { name: 'All Supported',
             extensions: ['cdxml','mol','sdf','xyz','pdb','cif'] }
         ] })
         -- Tauri dialog plugin shows native OS file picker
      |
      v  readTextFile(filePath)
         -- Tauri fs plugin reads the file as a UTF-8 string
      |
      v  Extension detection:

         .cdxml -> cdxmlToState(content)
                   [src/utils/cdxml.ts -- pure TypeScript XML parser]
                   Returns CanvasState directly with atoms, bonds, AND arrows
                   -> canvasRef.current.loadNative(content)

         .mol / .sdf / .xyz / .pdb / .cif
               -> invoke('parse_chemical_file', { content, format: ext })
                  -> Rust IPC -> Python sidecar `parse` subcommand
                  -> RDKit: Chem.MolFromMolBlock() / MolFromXYZBlock() / etc.
                  -> If no 3D coords: AllChem.Compute2DCoords(mol)
                  -> Returns { molblock: "...", smiles: "..." }
               -> canvasRef.current.loadFromSmiles(result.molblock)
      |
      v  pushToHistory(newState)   -- commits to undo stack
         Canvas re-renders; 200ms debounce fires -> SMILES updated
\end{lstlisting}

\subsection{Saving a File}

\begin{lstlisting}[language={}]
User presses Ctrl+S (default: saves as .cdxml)
      |
      v  save({ filters: [{ name: 'ChemDraw XML', extensions: ['cdxml'] }] })
         -- Tauri dialog plugin shows native OS save dialog
      |
      v  canvasRef.current.saveNative()
         ->  stateToCDXML(atoms, bonds, arrows)  [src/utils/cdxml.ts]
         ->  writeTextFile(path, xml)  [Tauri fs plugin]

File -> Export As -> Molfile:
      v  canvasRef.current.getMolblock()
         -> graphToMolblock(atoms, bonds, width, height)
      v  invoke('export_chemical_file', { molblock, format: 'mol' })
         -> Python: Chem.MolFromMolBlock(molblock) -> Chem.MolToMolBlock(mol)
      v  writeTextFile(path, result.content)

RXN export (unique -- requires arrow data):
      v  invoke('export_chemical_file', {
           molblock, format: 'rxn',
           arrowsJson: JSON.stringify(arrows)
         })
         -> Python: Chem.GetMolFrags(mol) gets disconnected fragments
         -> Compares each fragment centroid to the reaction arrow midpoint
         -> Fragments left of midpoint -> reactants; right -> products
         -> Writes "$RXN\n...\n$MOL\n..." format

PNG export:
      v  canvasRef.current.exportPNG()
         -> stageRef.current.toDataURL()  [Konva: renders canvas to base64 PNG]
         -> writeFile(path, Uint8Array from base64)  [binary write]

SVG export:
      v  canvasRef.current.exportSVG()
         -> generateCanvasSVG(atoms, bonds, arrows, ...)  [src/lib/svgExport.ts]
         -> writeTextFile(path, svg)

3D viewer screenshot:
      v  molecule3DRef.current.captureViewer()
         -> Re-renders the Three.js scene at a higher temporary pixel ratio
         -> Returns data URL (base64 PNG)
      v  writeFile(path, Uint8Array from base64)
\end{lstlisting}

% ─────────────────────────────────────────────────────────────────────────────
\section{Adding Features: A Practical Guide}
\label{sec:adding}

\subsection{How to Add a New Tauri Command}

Suppose you want a command that computes molecular formula and weight from SMILES.

\textbf{Step 1 --- Python sidecar} (\texttt{src-tauri/bin/chem-engine.py}):
\begin{lstlisting}[language=Python]
def get_mol_info(smiles):
    from rdkit.Chem import rdMolDescriptors, Descriptors
    mol = Chem.MolFromSmiles(smiles)
    if not mol: return {"error": "Invalid SMILES"}
    return {
        "formula": rdMolDescriptors.CalcMolFormula(mol),
        "mw": round(Descriptors.ExactMolWt(mol), 4),
    }

# In the argparse section:
info_parser = subparsers.add_parser("info")
info_parser.add_argument("smiles")

# In the dispatch section:
elif args.command == "info":
    print(json.dumps(get_mol_info(args.smiles)))
\end{lstlisting}

\textbf{Step 2 --- Rust command} (\texttt{src-tauri/src/lib.rs}):
\begin{lstlisting}[language=Rust]
#[tauri::command]
async fn get_mol_info(app: tauri::AppHandle, smiles: String)
    -> Result<Value, String>
{
    run_engine(app, vec!["info".to_string(), smiles], None).await
}

// In the invoke_handler list:
.invoke_handler(tauri::generate_handler![
    generate_conformers, parse_chemical_file, export_chemical_file,
    get_mol_info   // <- add here
])
\end{lstlisting}

\textbf{Step 3 --- TypeScript} (wherever you need it):
\begin{lstlisting}[language=TypeScript]
import { invoke } from '@tauri-apps/api/core';

const result: any = await invoke('get_mol_info', { smiles: 'CCO' });
if (!result.error) {
  console.log(result.formula);  // "C2H6O"
  console.log(result.mw);       // 46.0419
}
\end{lstlisting}

That is the complete cycle. No HTTP routes, no serialisation boilerplate.

\subsection{How to Add a New Drawing Tool}

\textbf{Step 1} --- Add the tool name to the \lstinline|Tool| union type in
\texttt{src/store/index.ts}:
\begin{lstlisting}[language=TypeScript]
export type Tool = 'select' | 'atom' | 'bond' | 'eraser' | 'fragment'
                 | 'arrow' | 'charge' | 'pan' | 'myTool';
\end{lstlisting}

\textbf{Step 2} --- Create a custom hook in \texttt{src/tools/useMyTool.ts}.
Follow the dual state+ref pattern for any values your callbacks need to read:

\begin{lstlisting}[language=TypeScript]
export function useMyTool() {
  // Any values read by onMouseMove/onMouseUp must be mirrored in refs:
  const [active, setActive] = useState(false);
  const activeRef = useRef(false);
  const setActiveBoth = (v: boolean) => { activeRef.current = v; setActive(v); };

  const onMouseDown = useCallback((pos: {x:number; y:number}, hitAtom: any) => {
    const { atoms, bonds, arrows, pushToHistory } = useStore.getState();
    // ... modify graph, push to history ...
    setActiveBoth(true);
  }, []);

  const onMouseMove = useCallback((pos: {x:number; y:number}) => {
    if (!activeRef.current) return;   // read ref, not state
    // ... update preview using setCurrentCanvasState (no history) ...
  }, []);

  const onMouseUp = useCallback(() => {
    if (!activeRef.current) return;
    setActiveBoth(false);
    const { atoms, bonds, arrows, pushToHistory } = useStore.getState();
    pushToHistory({ atoms, bonds, arrows });  // commit the action
  }, []);

  const cleanup = useCallback(() => { setActiveBoth(false); }, []);

  return { active, onMouseDown, onMouseMove, onMouseUp, cleanup };
}
\end{lstlisting}

\textbf{Step 3} --- Instantiate and dispatch in \texttt{ChemCanvas.tsx}:
\begin{lstlisting}[language=TypeScript]
const myTool = useMyTool();

// In the cleanup effect:
if (prevToolRef.current !== tool) { myTool.cleanup(); /* ... */ }

// In handleMouseDown:
if (tool === 'myTool') myTool.onMouseDown(pos, hitAtom);
\end{lstlisting}

\textbf{Step 4} --- Add a toolbar button in \texttt{App.tsx}:
\begin{lstlisting}[language=TypeScript]
<div onClick={() => setTool('myTool')} style={toolBtnStyle(tool === 'myTool')}>
  {/* icon */}
</div>
\end{lstlisting}

\subsection{How to Add a New Atom Shorthand}

In \texttt{src/lib/graph.ts}, add an entry to \lstinline|SHORTHAND_DATA|:
\begin{lstlisting}[language=TypeScript]
export const SHORTHAND_DATA: Record<string, { lead: string; valency: number }> = {
  // ...existing entries...
  'Ts':  { lead: 'S', valency: 2 },   // tosyl -- lead atom is S
  'Boc': { lead: 'O', valency: 2 },   // tert-butyloxycarbonyl -- lead is O
};
\end{lstlisting}

\lstinline|lead| is the element written to the V2000 Molblock. \lstinline|valency| is
the number of bonds the shorthand group can form. To also add a hotkey, add to
\lstinline|ATOM_HOTKEYS| in \texttt{ChemCanvas.tsx}:
\begin{lstlisting}[language=TypeScript]
const ATOM_HOTKEYS: Record<string, string> = {
  // ...existing...
  't': 'Ts',   // pressing 't' while hovering an atom switches it to Ts
};
\end{lstlisting}

\subsection{How to Add a New Ring Template}

In \texttt{App.tsx}, add to the \lstinline|FRAGMENTS| array:
\begin{lstlisting}[language=TypeScript]
const FRAGMENTS = [
  // ...existing...
  { id: 'indane', label: 'ind', smiles: 'C1CCc2ccccc21' },
];
\end{lstlisting}

The SMILES string is passed to RDKit for coordinate generation.
\lstinline|placeFragmentAt| handles placement automatically. To also add a keyboard
shortcut, add to \lstinline|BOND_HOTKEYS| in \texttt{ChemCanvas.tsx}:
\begin{lstlisting}[language=TypeScript]
'i': { ring: 'C1CCc2ccccc21' },   // 'i' key places indane
\end{lstlisting}

\subsection{How to Add a New File Format}

\textbf{For parsing (opening)}:
\begin{enumerate}[leftmargin=1.5em]
  \item Add the extension to the \lstinline|open()| dialog filter in \texttt{App.tsx}.
  \item Add it to the \lstinline|['mol', 'sdf', 'xyz', ...]| branch.
  \item In \texttt{chem-engine.py}, add a new branch to \lstinline|parse_file()|.
\end{enumerate}

\textbf{For exporting (saving)}:
\begin{enumerate}[leftmargin=1.5em]
  \item Add the format to the export submenu in \texttt{App.tsx}.
  \item Add a handler in \lstinline|handleExport|.
  \item In \texttt{chem-engine.py}, add a new branch to \lstinline|export_file()|:
\end{enumerate}
\begin{lstlisting}[language=Python]
elif format == 'inchi':
    from rdkit.Chem.inchi import MolToInchi
    return {"content": MolToInchi(mol)}
\end{lstlisting}

If the format can be handled entirely in TypeScript, skip the Python steps and write
a converter in \texttt{src/utils/}.

% ─────────────────────────────────────────────────────────────────────────────
\section{Critical Files Reference}
\label{sec:files}

\begin{table}[h!]
\centering
\begin{tabularx}{\textwidth}{@{} l X @{}}
\toprule
\textbf{File} & \textbf{Purpose} \\
\midrule
\texttt{src/types/chemistry.ts} & Data model: \texttt{Atom}, \texttt{Bond}, \texttt{Arrow}, \texttt{CanvasState} \\
\texttt{src/store/index.ts} & Zustand store: all global state and actions \\
\texttt{src/lib/graph.ts} & \texttt{graphToMolblock}, \texttt{molblockToState}, \texttt{getConnectedComponents}, \texttt{hashCanvasState} \\
\texttt{src/lib/svgExport.ts} & Canvas $\to$ SVG vector export \\
\texttt{src/hooks/useRDKit.ts} & Singleton WASM loader \\
\texttt{src/tools/useBondTool.ts} & Bond drawing logic (dual state+ref pattern) \\
\texttt{src/tools/useSelectTool.ts} & Selection, drag-move, rubber-band, rotation \\
\texttt{src/tools/useAtomTool.ts} & Atom placement and inline label editing \\
\texttt{src/components/ChemCanvas.tsx} & Drawing engine: Konva canvas, hit testing, zoom/pan, SMILES pipeline \\
\texttt{src/components/Molecule.tsx} & 2D SVG preview via RDKit WASM + Blob URL \\
\texttt{src/components/Molecule3DThreePanel.tsx} & 3D viewer shell: controls, conformer strip, ref bridge to the renderer \\
\texttt{src/components/Molecule3DThree.tsx} & Three.js renderer: scene, lights, meshes, measurement interaction, PNG capture \\
\texttt{src/components/IrPanel.tsx} & IR spectrum panel: xtb prediction, Lorentzian broadening, SVG chart \\
\texttt{src/components/MsPanel.tsx} & Mass spectrum panel: RDKit isotope pattern, SVG chart \\
\texttt{src/components/NmrPanel.tsx} & NMR panel: custom HOSE/DFT viewer, reversed-axis SVG chart, peak/atom hover linkage \\
\texttt{src/utils/cdxml.ts} & ChemDraw CDXML read/write \\
\texttt{src/App.tsx} & Root component: layout, theme, menus, keyboard shortcuts, viewer modes \\
\texttt{src-tauri/src/lib.rs} & Rust command handlers and sidecar runner \\
\texttt{src-tauri/src/main.rs} & Rust entry point (\texttt{windows\_subsystem} attribute) \\
\texttt{src-tauri/bin/chem-engine.py} & Python sidecar: conformers (ETKDGv3), IR (xtb), NMR (HOSE + optional DFT), parse, export \\
\texttt{src-tauri/bin/nmr\_hose.py} & HOSE code generation, lookup, and fast proton/carbon heuristics \\
\texttt{src-tauri/bin/nmr\_hose\_db.json} & Bundled HOSE lookup database generated from NMRShiftDB2 \\
\texttt{src-tauri/build-sidecar.js} & PyInstaller build; downloads xtb, builds HOSE DB, handles triple-naming \\
\texttt{development/build\_nmr\_hose\_db.py} & Offline HOSE DB build from NMRShiftDB2 \\
\texttt{src-tauri/tauri.conf.json} & Window size, decorations, CSP, \texttt{externalBin} \\
\texttt{vite.config.ts} & Dev server port 1420, HMR, ignored watch paths \\
\texttt{package.json} & npm scripts (including predev RDKit copy), JS dependencies \\
\texttt{src-tauri/Cargo.toml} & Rust dependencies (\texttt{tauri}, \texttt{serde\_json}, \texttt{tauri-plugin-*}) \\
\bottomrule
\end{tabularx}
\caption{Critical files and their roles}
\end{table}

\end{document}
